\documentclass[aspectratio=169]{beamer}
\usetheme{metropolis}           % Tema moderno y limpio
\useoutertheme{metropolis}
\useinnertheme{metropolis}
\usecolortheme{default}
\usepackage[spanish]{babel}
\usepackage[utf8]{inputenc}
\usepackage{amsmath, amssymb, amsthm}
\usepackage{mathrsfs}
\usepackage{mathtools}
\usepackage{graphicx}
\usepackage{dsfont}
\usepackage{tikz}
\usepackage{pgfplots}
\pgfplotsset{compat=1.17}
\usetikzlibrary{automata, positioning, arrows.meta, shapes.geometric}

% Configuración de colores
\setbeamercolor{title}{fg=black!90!orange}
\setbeamercolor{frametitle}{fg=white, bg=orange!70!black}
\setbeamercolor{block title}{fg=white, bg=orange!60!black}
\setbeamercolor{block body}{fg=black, bg=gray!5}
\setbeamercolor{example block title}{fg=white, bg=green!50!black}
\setbeamercolor{example block body}{fg=black, bg=green!5}

% Definir entornos de teoremas
\newtheorem{teorema}{Teorema}
\newtheorem{definicion}{Definición}

% Comandos útiles
\newcommand{\R}{\mathbb{R}}
\newcommand{\E}{\mathbb{E}}
\newcommand{\Prob}{\mathbb{P}}
\newcommand{\norm}[1]{\left\lVert #1 \right\rVert}
\newcommand{\abs}[1]{\left\lvert #1 \right\rvert}
\newcommand{\paren}[1]{\left( #1 \right)}
\newcommand{\bracks}[1]{\left[ #1 \right]}
\newcommand{\argmin}{\operatorname{argmin}}
\newcommand{\argmax}{\operatorname{argmax}}

% Información de la presentación
\title{Algoritmo K-means}
\subtitle{Un recorrido por la formulación matemática, convergencia y variantes}
\author{}
\date{}

\begin{document}

\begin{frame}
    \titlepage
\end{frame}

\begin{frame}{Contenido}
    \tableofcontents
\end{frame}

\section{Introducción al clustering}

\begin{frame}{¿Qué es el clustering?}
    \begin{itemize}
        \item El \textbf{clustering} (o agrupamiento) es una técnica de aprendizaje no supervisado.
        \item Objetivo: particionar un conjunto de datos en grupos (clústeres) tales que:
        \begin{itemize}
            \item Los puntos dentro de un mismo clúster sean \textbf{similares} entre sí.
            \item Los puntos de diferentes clústeres sean \textbf{disímiles}.
        \end{itemize}
        \item Aplicaciones:
        \begin{itemize}
            \item Segmentación de clientes
            \item Compresión de imágenes
            \item Análisis de datos biológicos
            \item Sistemas de recomendación
        \end{itemize}
    \end{itemize}
    
    \begin{center}
        \begin{tikzpicture}[scale=0.6]
            % Datos sin agrupar
            \begin{scope}[xshift=0cm]
                \draw[gray!30] (0,0) rectangle (4,4);
                \foreach \x/\y in {1/2, 1.2/1.8, 0.8/2.2, 1.5/3, 1.8/2.5, 2.5/1.2, 2.8/1, 3/1.5, 3.2/0.8, 2.2/1.8, 3.5/2.8, 3.2/3, 3.8/2.5, 3.5/2, 2.8/3.2} {
                    \fill[blue] (\x,\y) circle (2pt);
                }
                \node at (2,4.3) {Datos originales};
            \end{scope}
            
            % Datos agrupados
            \begin{scope}[xshift=5cm]
                \draw[gray!30] (0,0) rectangle (4,4);
                \foreach \x/\y in {1/2, 1.2/1.8, 0.8/2.2, 1.5/3, 1.8/2.5, 2.2/1.8} {
                    \fill[red] (\x,\y) circle (2pt);
                }
                \foreach \x/\y in {2.5/1.2, 2.8/1, 3/1.5, 3.2/0.8} {
                    \fill[blue] (\x,\y) circle (2pt);
                }
                \foreach \x/\y in {3.5/2.8, 3.2/3, 3.8/2.5, 3.5/2, 2.8/3.2} {
                    \fill[green!70!black] (\x,\y) circle (2pt);
                }
                \node at (2,4.3) {Clustering (\(K=3\))};
            \end{scope}
        \end{tikzpicture}
    \end{center}
\end{frame}

\begin{frame}{Tipos de algoritmos de clustering}
    \begin{columns}
        \column{0.5\textwidth}
        \textbf{Jerárquicos:}
        \begin{itemize}
            \item Construyen una jerarquía de clústeres.
            \item Dendrograma como resultado.
            \item Ejemplos: single-link, complete-link.
        \end{itemize}
        
        \column{0.5\textwidth}
        \textbf{Particionales:}
        \begin{itemize}
            \item Dividen los datos en \(K\) particiones.
            \item Requieren especificar \(K\) de antemano.
            \item Ejemplos: \textbf{K-means}, K-medoids.
        \end{itemize}
    \end{columns}
    
    \vspace{0.5cm}
    \begin{center}
        \textbf{K-means es el algoritmo particional más popular por su simplicidad y eficiencia.}
    \end{center}
\end{frame}

\section{Formulación matemática del problema}

\begin{frame}{Definición formal}
    \begin{block}{Datos de entrada}
        \begin{itemize}
            \item Conjunto de datos: \(X = \{x_1, x_2, \ldots, x_n\} \subset \R^d\)
            \item Número de clústeres: \(K \in \mathbb{N}\) (fijado de antemano)
        \end{itemize}
    \end{block}
    
    \begin{block}{Objetivo}
        Encontrar una partición de \(X\) en \(K\) subconjuntos \(C_1, C_2, \ldots, C_K\) y sus respectivos centroides \(\mu_1, \mu_2, \ldots, \mu_K \in \R^d\) que minimicen la \textbf{función de error}:
        
        \[
        J(C, \mu) = \sum_{j=1}^{K} \sum_{x_i \in C_j} \norm{x_i - \mu_j}^2
        \]
    \end{block}
    
    \pause
    \textbf{Interpretación:} \(J\) mide la varianza intra-clúster (inertia). Queremos que los puntos estén lo más cerca posible de su centroide.
\end{frame}

\begin{frame}{Notación y terminología}
    \begin{itemize}
        \item \textbf{Centroide \(\mu_j\)}: representante del clúster \(j\), calculado como la media de sus puntos.
        \item \textbf{Distancia euclídea}: \(\norm{x_i - \mu_j}^2 = \sum_{m=1}^d (x_{im} - \mu_{jm})^2\)
        \item \textbf{Participación}: cada punto pertenece a \textbf{un único} clúster (hard clustering).
        \item \textbf{Función de costo}: también llamada \textbf{inertia} o \textbf{SSE} (Sum of Squared Errors).
    \end{itemize}
    
    \begin{exampleblock}{Ejemplo unidimensional}
        \(X = \{1, 2, 3, 8, 9, 10\}\), \(K=2\)
        
        Si \(C_1 = \{1,2,3\}\), \(\mu_1 = 2\); \(C_2 = \{8,9,10\}\), \(\mu_2 = 9\)
        
        \(J = (1-2)^2 + (2-2)^2 + (3-2)^2 + (8-9)^2 + (9-9)^2 + (10-9)^2 = 1+0+1+1+0+1 = 4\)
    \end{exampleblock}
\end{frame}

\begin{frame}{El problema de optimización}
    \begin{block}{Problema de optimización combinatoria}
        \[
        \min_{\substack{\text{particiones }\{C_j\}_{j=1}^K \\ \text{centroides } \mu_j \in \R^d}} \sum_{j=1}^K \sum_{x_i \in C_j} \norm{x_i - \mu_j}^2
        \]
    \end{block}
    
    \pause
    \textbf{Dificultades:}
    \begin{itemize}
        \item Problema \textbf{NP-hard} en general (para \(K>1\) y \(d>1\)).
        \item Número de particiones posibles: \(\frac{1}{K!} \sum_{j=1}^K (-1)^{K-j} \binom{K}{j} j^n\) (crecimiento exponencial).
        \item Función no convexa y no diferenciable por la asignación discreta.
    \end{itemize}
    
    \pause
    \textbf{Solución práctica:} Algoritmos iterativos que encuentran \textbf{mínimos locales}.
\end{frame}

\section{Algoritmo de Lloyd (K-means clásico)}

\begin{frame}{Historia}
    \begin{itemize}
        \item El algoritmo fue propuesto por \textbf{Stuart Lloyd} en 1957 en los Laboratorios Bell, pero no se publicó hasta 1982 [Lloyd, 1982].
        \item Independientemente, \textbf{Hugo Steinhaus} (1956) y \textbf{Edward Forgy} (1965) propusieron variantes similares.
        \item Por eso a veces se llama \textbf{algoritmo de Lloyd-Forgy}.
        \item Es uno de los algoritmos más utilizados en minería de datos.
    \end{itemize}
    
    \begin{block}{Referencia clásica}
        Lloyd, S. P. (1982). Least squares quantization in PCM. \textit{IEEE Transactions on Information Theory}, 28(2), 129-137.
    \end{block}
\end{frame}

\begin{frame}{Estructura general del algoritmo}
    \begin{block}{Algoritmo de K-means (Lloyd)}
        \textbf{Entrada:} Datos \(X = \{x_1,\ldots,x_n\}\), número de clústeres \(K\), criterio de parada.
        
        \textbf{Inicialización:} Elegir \(K\) centroides iniciales \(\mu_1^{(0)}, \ldots, \mu_K^{(0)}\).
        
        \textbf{Iteración para \(t = 0, 1, 2, \ldots\):}
        \begin{enumerate}
            \item \textbf{Asignación:} Para cada punto \(x_i\), asignarlo al centroide más cercano:
            \[
            C_j^{(t)} = \left\{ x_i : \norm{x_i - \mu_j^{(t)}}^2 \leq \norm{x_i - \mu_\ell^{(t)}}^2 \; \forall \ell \right\}
            \]
            (en caso de empate, asignar arbitrariamente).
            
            \item \textbf{Actualización:} Recalcular cada centroide como la media de los puntos asignados:
            \[
            \mu_j^{(t+1)} = \frac{1}{|C_j^{(t)}|} \sum_{x_i \in C_j^{(t)}} x_i
            \]
        \end{enumerate}
        
        Hasta que los centroides no cambien significativamente o se alcance el máximo de iteraciones.
    \end{block}
\end{frame}

\begin{frame}{Visualización del proceso iterativo}
    \begin{center}
        \begin{tikzpicture}[scale=0.7]
            % Iteración 1
            \begin{scope}[xshift=0cm]
                \draw[gray!30] (0,0) rectangle (4,4);
                \foreach \x/\y in {0.5/1, 1/0.5, 1.5/1.5, 2/1, 2.5/3, 3/2.5, 3.5/3.5} {
                    \fill[blue] (\x,\y) circle (2pt);
                }
                \fill[red] (1,1) circle (4pt) node[below right] {\(\mu_1\)};
                \fill[red] (3,3) circle (4pt) node[below right] {\(\mu_2\)};
                \node at (2,-0.5) {Inicialización};
            \end{scope}
            
            % Iteración 2
            \begin{scope}[xshift=5cm]
                \draw[gray!30] (0,0) rectangle (4,4);
                \foreach \x/\y in {0.5/1, 1/0.5, 1.5/1.5, 2/1} {
                    \fill[blue!50!red] (\x,\y) circle (2pt);
                }
                \foreach \x/\y in {2.5/3, 3/2.5, 3.5/3.5} {
                    \fill[blue!50!green] (\x,\y) circle (2pt);
                }
                \fill[red] (1.2,1) circle (4pt);
                \fill[red] (3,3) circle (4pt);
                \node at (2,-0.5) {Asignación};
            \end{scope}
            
            % Iteración 3
            \begin{scope}[xshift=10cm]
                \draw[gray!30] (0,0) rectangle (4,4);
                \foreach \x/\y in {0.5/1, 1/0.5, 1.5/1.5, 2/1} {
                    \fill[blue!50!red] (\x,\y) circle (2pt);
                }
                \foreach \x/\y in {2.5/3, 3/2.5, 3.5/3.5} {
                    \fill[blue!50!green] (\x,\y) circle (2pt);
                }
                \fill[red] (1.3,0.9) circle (4pt);
                \fill[red] (3.2,3) circle (4pt);
                \node at (2,-0.5) {Actualización};
            \end{scope}
            
            % Iteración final
            \begin{scope}[xshift=15cm]
                \draw[gray!30] (0,0) rectangle (4,4);
                \foreach \x/\y in {0.5/1, 1/0.5, 1.5/1.5, 2/1} {
                    \fill[blue!50!red] (\x,\y) circle (2pt);
                }
                \foreach \x/\y in {2.5/3, 3/2.5, 3.5/3.5} {
                    \fill[blue!50!green] (\x,\y) circle (2pt);
                }
                \fill[red] (1.25,1) circle (4pt);
                \fill[red] (3,2.9) circle (4pt);
                \node at (2,-0.5) {Convergencia};
            \end{scope}
        \end{tikzpicture}
    \end{center}
    
    \textbf{Evolución:} Los centroides se mueven hacia el centro de sus clusters asignados hasta estabilizarse.
\end{frame}

\section{Propiedades de convergencia}

\begin{frame}{Teorema de convergencia finita}
    \begin{teorema}[Selim \& Ismail, 1984]
        El algoritmo de K-means (Lloyd) converge en un \textbf{número finito de iteraciones} para cualquier métrica, bajo la condición de que no haya empates en la asignación.
    \end{teorema}
    
    \begin{block}{Idea de la demostración}
        \begin{enumerate}
            \item El número de particiones posibles de \(n\) puntos en \(K\) clusters es \textbf{finito}.
            \item Cada iteración (asignación + actualización) \textbf{reduce} estrictamente la función de costo \(J\), o el algoritmo se detiene.
            \item Si \(J\) no puede disminuir indefinidamente (está acotada inferiormente por 0), el algoritmo debe alcanzar un mínimo local tras un número finito de pasos.
        \end{enumerate}
    \end{block}
    
    \begin{alertblock}{Importante}
        El algoritmo converge a un \textbf{mínimo local}, no necesariamente al mínimo global. El resultado final depende de la inicialización.
    \end{alertblock}
\end{frame}

\begin{frame}{¿Por qué solo mínimo local?}
    \begin{itemize}
        \item La función \(J(C, \mu)\) no es convexa debido a la asignación discreta de puntos a clusters.
        \item El algoritmo es un caso especial de \textbf{descenso por coordenadas} (coordinate descent):
        \begin{itemize}
            \item Paso de asignación: minimiza \(J\) sobre las asignaciones manteniendo fijos los centroides.
            \item Paso de actualización: minimiza \(J\) sobre los centroides manteniendo fijas las asignaciones.
        \end{itemize}
        \item Cada paso reduce \(J\), pero el algoritmo puede quedar atrapado en un mínimo local.
    \end{itemize}
    
    \begin{center}
        \begin{tikzpicture}[scale=0.6]
            \begin{axis}[
                xlabel={Iteración},
                ylabel={\(J\)},
                xmin=0, xmax=10,
                ymin=0, ymax=100,
                width=8cm, height=5cm,
                grid=both
            ]
            \addplot[blue, thick, mark=*] coordinates {
                (0,95) (1,70) (2,52) (3,40) (4,33) (5,28) (6,25) (7,23) (8,22) (9,21.5) (10,21.5)
            };
            \end{axis}
        \end{tikzpicture}
    \end{center}
\end{frame}

\begin{frame}{Condiciones de optimalidad}
    \begin{definicion}[Punto de Kuhn-Tucker]
        Una solución \((\{C_j\}, \{\mu_j\})\) es un punto de Kuhn-Tucker (condición necesaria de optimalidad) si:
        \begin{enumerate}
            \item Para cada \(j\), \(\mu_j = \frac{1}{|C_j|} \sum_{x_i \in C_j} x_i\) (condición de centroide).
            \item Para cada \(x_i \in C_j\), \(\norm{x_i - \mu_j} \leq \norm{x_i - \mu_\ell}\) para todo \(\ell \neq j\) (condición de asignación óptima).
        \end{enumerate}
    \end{definicion}
    
    \begin{teorema}[Selim \& Ismail, 1984]
        Bajo condiciones de diferenciabilidad, el algoritmo de K-means converge a un punto de Kuhn-Tucker. Si además la función es convexa en cada región de asignación, se alcanza un mínimo local.
    \end{teorema}
\end{frame}

\section{Inicialización y sensibilidad}

\begin{frame}{Problema de la inicialización}
    \begin{center}
        \textbf{El algoritmo es muy sensible a la elección inicial de los centroides.}
    \end{center}
    
    \begin{columns}
        \column{0.5\textwidth}
        \textbf{Mala inicialización:}
        \begin{itemize}
            \item Centroides lejos de los datos
            \item Clusters vacíos
            \item Mínimo local pobre
        \end{itemize}
        
        \column{0.5\textwidth}
        \textbf{Buena inicialización:}
        \begin{itemize}
            \item Centroides cerca de los datos
            \item Convergencia rápida
            \item Mejor solución
        \end{itemize}
    \end{columns}
    
    \vspace{0.3cm}
    \begin{exampleblock}{Ejemplo}
        Con \(K=3\) y datos con 3 grupos naturales, si dos centroides iniciales caen en el mismo grupo, ese grupo se partirá artificialmente y otro grupo quedará sin representante.
    \end{exampleblock}
\end{frame}

\begin{frame}{Métodos de inicialización comunes}
    \begin{itemize}
        \item \textbf{Aleatoria simple:} Elegir \(K\) puntos al azar del conjunto de datos.
        \begin{itemize}
            \item Simple pero puede dar malos resultados.
        \end{itemize}
        
        \item \textbf{Método de Forgy:} Elegir \(K\) observaciones al azar como centroides iniciales.
        \begin{itemize}
            \item Similar a la aleatoria simple.
        \end{itemize}
        
        \item \textbf{Inicialización por partición aleatoria:} Asignar cada punto a un cluster aleatoriamente, luego calcular medias.
        
        \item \textbf{K-means++:} Inicialización probabilística que maximiza la distancia entre centroides [Arthur \& Vassilvitskii, 2007].
        \begin{itemize}
            \item Primero elige un centroide al azar.
            \item Cada nuevo centroide se elige con probabilidad proporcional al cuadrado de la distancia al centroide más cercano ya elegido.
            \item Garantiza una solución \(O(\log K)\)-competitiva.
        \end{itemize}
    \end{itemize}
\end{frame}

\begin{frame}{Algoritmo K-means++}
    \begin{block}{Inicialización K-means++}
        \textbf{Entrada:} Datos \(X\), número de clusters \(K\).
        
        \begin{enumerate}
            \item Elegir el primer centroide \(\mu_1\) uniformemente al azar de \(X\).
            \item Para \(j = 2\) hasta \(K\):
            \begin{itemize}
                \item Para cada punto \(x\), calcular \(D(x) = \min_{\ell < j} \norm{x - \mu_\ell}^2\).
                \item Elegir el siguiente centroide \(\mu_j\) de \(X\) con probabilidad \(\frac{D(x)}{\sum_{x' \in X} D(x')}\).
            \end{itemize}
        \end{enumerate}
        
        Luego ejecutar K-means estándar con estos centroides iniciales.
    \end{block}
    
    \textbf{Ventaja:} Mejora significativamente la calidad y la consistencia de los resultados.
\end{frame}

\section{Elección del número de clusters \(K\)}

\begin{frame}{¿Cómo elegir \(K\)?}
    \begin{itemize}
        \item En aplicaciones reales, normalmente \textbf{no sabemos} el número óptimo de clusters.
        \item Necesitamos métodos para determinar un \(K\) adecuado.
        \item Compromiso:
        \begin{itemize}
            \item \(K\) muy pequeño: sub-agrupamiento (underfitting).
            \item \(K\) muy grande: sobre-agrupamiento (overfitting).
        \end{itemize}
    \end{itemize}
    
    \begin{center}
        \textbf{Buscamos un equilibrio entre simplicidad y fidelidad a los datos.}
    \end{center}
\end{frame}

\begin{frame}{Método del codo (Elbow Method)}
    \begin{columns}
        \column{0.5\textwidth}
        \begin{enumerate}
            \item Ejecutar K-means para \(K = 1, 2, \ldots, K_{max}\).
            \item Calcular la inercia \(J(K)\) para cada \(K\).
            \item Graficar \(J(K)\) vs \(K\).
            \item Buscar el "codo": punto donde la mejora empieza a ser marginal.
        \end{enumerate}
        
        \column{0.5\textwidth}
        \begin{tikzpicture}[scale=0.8]
            \begin{axis}[
                xlabel={\(K\)},
                ylabel={Inercia \(J(K)\)},
                xmin=1, xmax=8,
                ymin=0, ymax=100,
                xtick={1,2,3,4,5,6,7,8},
                width=6cm, height=5cm,
                grid=both
            ]
            \addplot[blue, thick, mark=*] coordinates {
                (1,95) (2,60) (3,35) (4,22) (5,18) (6,15) (7,13) (8,11.5)
            };
            \draw[red, dashed, thick] (axis cs:3,35) -- (axis cs:3,0);
            \node[red] at (axis cs:3.5,40) {Codo};
            \end{axis}
        \end{tikzpicture}
    \end{columns}
    
    \begin{itemize}
        \item El "codo" indica el \(K\) donde la ganancia de añadir más clusters deja de ser significativa.
        \item Subjetivo: a veces no hay un codo claro.
    \end{itemize}
\end{frame}

\begin{frame}{Coeficiente de silueta}
    \begin{definicion}
        Para un punto \(x_i\) en el cluster \(C_j\):
        \begin{itemize}
            \item \(a(i) =\) distancia media de \(x_i\) a los demás puntos de su mismo cluster (cohesión).
            \item \(b(i) =\) distancia media de \(x_i\) a los puntos del cluster vecino más cercano (separación).
        \end{itemize}
        Coeficiente de silueta:
        \[
        s(i) = \frac{b(i) - a(i)}{\max\{a(i), b(i)\}} \in [-1, 1]
        \]
    \end{definicion}
    
    \textbf{Interpretación:}
    \begin{itemize}
        \item \(s(i) \approx 1\): punto bien clasificado.
        \item \(s(i) \approx 0\): punto en la frontera entre dos clusters.
        \item \(s(i) < 0\): punto probablemente mal clasificado.
    \end{itemize}
    
    \textbf{Métrica global:} Silueta promedio para todos los puntos. Elegir \(K\) que maximice la silueta promedio.
\end{frame}

\begin{frame}{Comparación de métodos para elegir \(K\)}
    \begin{table}
        \centering
        \begin{tabular}{|l|c|c|}
            \hline
            \textbf{Método} & \textbf{Ventajas} & \textbf{Desventajas} \\
            \hline
            Método del codo & Simple, intuitivo & Subjetivo \\
            Silueta & Objetivo, interpretable & Costoso computacionalmente \\
            Estadístico Gap & Riguroso estadísticamente & Complejo \\
            Regla de la raíz de \(n/2\) & Muy simple & Heurística sin garantías \\
            \hline
        \end{tabular}
    \end{table}
    
    \textbf{Recomendación:} Usar múltiples métodos y validación cruzada cuando sea posible.
\end{frame}

\section{Complejidad computacional}

\begin{frame}{Análisis de complejidad}
    \begin{block}{Por iteración}
        \begin{itemize}
            \item \textbf{Asignación:} Calcular distancia de cada punto a cada centroide: \(O(n \cdot K \cdot d)\)
            \begin{itemize}
                \item \(n\): número de puntos
                \item \(K\): número de clusters
                \item \(d\): dimensionalidad
            \end{itemize}
            \item \textbf{Actualización:} Recalcular medias: \(O(n \cdot d)\)
        \end{itemize}
    \end{block}
    
    \begin{block}{Complejidad total}
        \(O(t \cdot n \cdot K \cdot d)\), donde \(t\) es el número de iteraciones hasta convergencia.
    \end{block}
    
    \begin{itemize}
        \item En la práctica, \(t\) suele ser pequeño (10-100 iteraciones).
        \item El algoritmo escala linealmente con \(n\), lo que lo hace adecuado para grandes conjuntos de datos.
        \item Para datos muy grandes, existen variantes como \textbf{mini-batch K-means} que usan submuestras.
    \end{itemize}
\end{frame}

\begin{frame}{Comparación con otros algoritmos}
    \begin{table}
        \centering
        \begin{tabular}{|l|c|c|c|}
            \hline
            \textbf{Algoritmo} & \textbf{Complejidad} & \textbf{Ventajas} & \textbf{Desventajas} \\
            \hline
            K-means & \(O(t n K d)\) & Rápido, simple & Asume clusters esféricos \\
            K-medoids & \(O(t (n-K)^2 K d)\) & Robusto a outliers & Lento \\
            DBSCAN & \(O(n \log n)\) & No requiere \(K\) & Sensible a parámetros \\
            Jerárquico & \(O(n^3)\) & Dendrograma & Muy lento \\
            \hline
        \end{tabular}
    \end{table}
    
    \textbf{K-means es el más eficiente para grandes volúmenes de datos cuando los clusters son aproximadamente esféricos.}
\end{frame}

\section{Variantes y extensiones}

\begin{frame}{Mini-batch K-means}
    \begin{block}{Idea}
        En lugar de usar todos los datos en cada iteración, usar pequeños subconjuntos aleatorios (mini-batches).
    \end{block}
    
    \textbf{Algoritmo:}
    \begin{enumerate}
        \item Inicializar centroides (ej. con K-means++).
        \item Repetir hasta convergencia:
        \begin{itemize}
            \item Seleccionar un mini-batch \(M\) de tamaño \(b \ll n\).
            \item Asignar cada punto de \(M\) al centroide más cercano.
            \item Actualizar centroides usando solo los puntos de \(M\) con una tasa de aprendizaje.
        \end{itemize}
    \end{enumerate}
    
    \begin{teorema}[Tang \& Monteleoni, 2017]
        Mini-batch K-means converge a una solución local óptima con tasa \(O(1/t)\) bajo condiciones generales.
    \end{teorema}
    
    \textbf{Ventaja:} Mucho más rápido para grandes conjuntos de datos.
\end{frame}

\begin{frame}{Online K-means (aprendizaje en línea)}
    \begin{block}{Configuración}
        Los datos llegan en flujo (streaming), uno a uno.
    \end{block}
    
    \begin{block}{Actualización para cada nuevo punto \(x\)}
        \begin{enumerate}
            \item Encontrar el centroide más cercano: \(j = \argmin_\ell \norm{x - \mu_\ell}\)
            \item Actualizar ese centroide:
            \[
            \mu_j \leftarrow \mu_j - \eta_j (\mu_j - x)
            \]
            donde \(\eta_j\) es una tasa de aprendizaje (p.ej., \(1/N_j\), siendo \(N_j\) el número de veces que se ha actualizado \(\mu_j\)).
        \end{enumerate}
    \end{block}
    
    \begin{teorema}[Dasgupta et al., 2022]
        El online K-means puede interpretarse como un descenso por gradiente estocástico y converge asintóticamente al conjunto de puntos estacionarios de la función de costo.
    \end{teorema}
\end{frame}

\begin{frame}{Kernel K-means}
    \begin{itemize}
        \item K-means estándar asume que los clusters son linealmente separables en el espacio original.
        \item \textbf{Kernel K-means} proyecta los datos a un espacio de características de mayor dimensión mediante una función kernel \(\phi\).
    \end{itemize}
    
    \begin{block}{Objetivo en espacio kernel}
        \[
        J = \sum_{j=1}^K \sum_{x_i \in C_j} \norm{\phi(x_i) - \mu_j^\phi}^2
        \]
        donde \(\mu_j^\phi = \frac{1}{|C_j|} \sum_{x_i \in C_j} \phi(x_i)\).
    \end{block}
    
    \begin{itemize}
        \item Usando el truco del kernel, podemos calcular distancias sin necesidad de conocer \(\phi\) explícitamente.
        \item Permite detectar clusters con formas no lineales.
    \end{itemize}
\end{frame}

\begin{frame}{K-medoids y K-medians}
    \begin{description}
        \item[K-medoids:] En lugar de la media, usa un punto real del conjunto como centroide (el \textit{medoide}). Más robusto a outliers.
        
        \item[K-medians:] Minimiza la suma de distancias \(L_1\) (mediana) en lugar de \(L_2\). Útil cuando los datos tienen colas pesadas.
    \end{description}
    
    \begin{table}
        \centering
        \begin{tabular}{|l|c|c|}
            \hline
            \textbf{Métrica} & \textbf{K-means} & \textbf{K-medians} \\
            \hline
            Distancia & Euclídea (\(L_2\)) & Manhattan (\(L_1\)) \\
            Centroide & Media & Mediana \\
            Robustez & Baja & Alta \\
            \hline
        \end{tabular}
    \end{table}
\end{frame}

\section{Aplicaciones y ejemplos}

\begin{frame}{Aplicación 1: Segmentación de clientes}
    \begin{columns}
        \column{0.5\textwidth}
        \textbf{Problema:} Una empresa quiere agrupar a sus clientes según:
        \begin{itemize}
            \item Ingreso anual
            \item Gasto anual
            \item Frecuencia de compra
        \end{itemize}
        
        \textbf{Objetivo:} Diseñar estrategias de marketing específicas para cada segmento.
        
        \column{0.5\textwidth}
        \begin{tikzpicture}[scale=0.7]
            \begin{axis}[
                xlabel={Ingreso},
                ylabel={Gasto},
                width=5cm, height=5cm,
                grid=both
            ]
            \addplot[only marks, mark=*, blue] coordinates {
                (20,5) (25,8) (30,6) (22,7) (28,10) (35,15) (40,18) (38,20) (42,22)
                (50,30) (55,35) (60,32) (58,38) (52,33)
            };
            \end{axis}
        \end{tikzpicture}
    \end{columns}
    
    \textbf{Resultado:} Segmentos como "clientes jóvenes con alto gasto", "clientes mayores conservadores", etc.
\end{frame}

\begin{frame}{Aplicación 2: Compresión de imágenes}
    \begin{itemize}
        \item Cada píxel es un punto en el espacio RGB (3 dimensiones).
        \item Aplicar K-means con \(K\) pequeño (ej. 16, 32, 64 colores).
        \item Reemplazar cada píxel por el color de su centroide más cercano.
    \end{itemize}
    
    \begin{center}
        \begin{tikzpicture}
            \draw[fill=red!80] (0,0) rectangle (2,2);
            \draw[fill=red!60] (2,0) rectangle (4,2);
            \draw[fill=red!40] (4,0) rectangle (6,2);
            \node at (1,-0.5) {Original};
            \node at (3,-0.5) {\(K=16\)};
            \node at (5,-0.5) {\(K=4\)};
        \end{tikzpicture}
    \end{center}
    
    \textbf{Resultado:} Reducción del tamaño del archivo manteniendo calidad aceptable.
\end{frame}

\begin{frame}{Aplicación 3: Detección de anomalías}
    \begin{itemize}
        \item Entrenar K-means con datos normales.
        \item Para un nuevo punto, medir su distancia al centroide más cercano.
        \item Si la distancia supera un umbral, marcar como anomalía.
    \end{itemize}
    
    \begin{center}
        \begin{tikzpicture}[scale=0.8]
            \draw[gray!30] (0,0) rectangle (5,4);
            \foreach \x/\y in {1/1, 1.2/1.3, 0.8/0.9, 1.5/1.8, 1.8/1.2, 2.2/2.5, 2.5/2.8, 2.3/2, 2.8/2.2, 3.2/3, 3.5/3.2, 3/2.8} {
                \fill[blue] (\x,\y) circle (2pt);
            }
            \fill[red] (4.5,3.5) circle (4pt);
            \draw[->, thick, red] (4.5,3.5) -- (2.5,2.2);
            \node[red] at (4.5,3.8) {Anomalía};
        \end{tikzpicture}
    \end{center}
    
    \textbf{Aplicaciones:} Detección de fraudes, intrusiones en redes, fallos en maquinaria.
\end{frame}

\section{Limitaciones y consideraciones}

\begin{frame}{Limitaciones de K-means}
    \begin{itemize}
        \item \textbf{Asume clusters esféricos:} Funciona mal con clusters alargados o con formas complejas.
        \item \textbf{Sensible a la escala:} Requiere normalización de características.
        \item \textbf{Sensible a outliers:} La media es afectada por valores extremos.
        \item \textbf{Requiere especificar \(K\):} No siempre es conocido.
        \item \textbf{Convergencia a mínimo local:} Depende de la inicialización.
        \item \textbf{Datos categóricos:} No funciona directamente (requiere transformación).
    \end{itemize}
\end{frame}

\begin{frame}{Cuándo usar K-means}
    \begin{block}{Casos favorables}
        \begin{itemize}
            \item Clusters aproximadamente esféricos y de tamaño similar.
            \item Datos numéricos en espacios de baja/media dimensión.
            \item Grandes volúmenes de datos (escalabilidad).
            \item Necesidad de interpretabilidad (centroides como representantes).
        \end{itemize}
    \end{block}
    
    \begin{block}{Casos desfavorables}
        \begin{itemize}
            \item Clusters con formas complejas o anidadas.
            \item Datos con outliers significativos.
            \item Datos categóricos sin codificación adecuada.
            \item Cuando se desconoce \(K\) y no hay método claro para determinarlo.
        \end{itemize}
    \end{block}
\end{frame}

\section{Resumen y conclusiones}

\begin{frame}{Resumen}
    \begin{itemize}
        \item \textbf{K-means} es un algoritmo de clustering particional que minimiza la suma de cuadrados intra-cluster.
        \item \textbf{Algoritmo de Lloyd:} iteraciones de asignación y actualización.
        \item \textbf{Convergencia:} finita a un mínimo local, no necesariamente global.
        \item \textbf{Sensible a:} inicialización (K-means++ ayuda) y elección de \(K\) (elbow, silueta).
        \item \textbf{Complejidad:} \(O(t n K d)\), escalable a grandes conjuntos.
        \item \textbf{Variantes:} Mini-batch, online, kernel, k-medoids.
        \item \textbf{Aplicaciones:} segmentación, compresión, detección de anomalías.
    \end{itemize}
\end{frame}

\begin{frame}{Referencias clave}
    \begin{thebibliography}{9}
        \bibitem{lloyd1982} Lloyd, S. P. (1982). Least squares quantization in PCM. \textit{IEEE Transactions on Information Theory}.
        
        \bibitem{selim1984} Selim, S. Z. \& Ismail, M. A. (1984). K-means-type algorithms: A generalized convergence theorem and characterization of local optimality. \textit{IEEE Transactions on Pattern Analysis and Machine Intelligence}.
        
        \bibitem{arthur2007} Arthur, D. \& Vassilvitskii, S. (2007). K-means++: The advantages of careful seeding. \textit{SODA}.
    \end{thebibliography}
\end{frame}

\begin{frame}{Referencias clave}
    \begin{thebibliography}{9}
        \bibitem{bottou1995} Bottou, L. \& Bengio, Y. (1995). Convergence properties of the K-means algorithms. \textit{NIPS}.
        
        \bibitem{dasgupta2022} Dasgupta, S., Mahajan, G. \& So, G. (2022). Convergence of online k-means. \textit{arXiv:2202.10640}.
        
        \bibitem{tang2017} Tang, C. \& Monteleoni, C. (2017). Convergence rate of stochastic k-means. \textit{arXiv:1611.05132}.
    \end{thebibliography}
\end{frame}

\end{document}